%! Absolute Value Equations \& Inequalities — Unit 1, Lesson 1.2 — Student Packet Cover Sheet
\documentclass[10pt]{article}
\usepackage{algebra2-article}
\usepackage{algebra2-boxes}

\begin{document}

% Full-bleed forest banner
\begin{tikzpicture}[remember picture, overlay]
  \fill[forest] (current page.north west)
    rectangle ([yshift=-0.9in]current page.north east);
\end{tikzpicture}
\vspace{-0.8in}

\begin{center}
  {\color{white}\textbf{\LARGE Algebra 2}}\\[4pt]
  {\color{forestmid}\large\textbf{Unit 1 \quad Linear Functions}}\\[6pt]
  {\color{white}\textbf{\large Lesson 1.2 \quad Absolute Value Equations \& Inequalities}}
\end{center}

\vspace{0.22in}
\namedateperiod
\vspace{0.18in}

\begin{learningtargetbox}
\small
\begin{enumerate}[leftmargin=1.5em, itemsep=5pt]
  \item \ldots{} read $|x-h|$ as the \textbf{distance} between $x$ and $h$, and explain why an
        absolute value equation usually has \emph{two} solutions.
  \item \ldots{} \textbf{isolate} the bars before splitting into two cases, and handle the special
        cases where there is exactly one solution or none at all.
  \item \ldots{} solve $|X| < c$ and $|X| > c$, and say which one gives a single interval
        (\emph{``less th\textbf{AND}''}) and which gives two rays (\emph{``great\textbf{OR}''}).
  \item \ldots{} write a solution set three ways --- \textbf{set notation}, \textbf{interval
        notation}, and a shaded \textbf{number line} --- and turn a ``within so much of''
        statement into math.
\end{enumerate}
\end{learningtargetbox}

\vspace{0.14in}

\begin{tocbox}
\small
\renewcommand{\arraystretch}{1.5}
\begin{tabularx}{\linewidth}{@{} c l X r @{}}
\rowcolor{forestbg}
\textbf{\#} & \textbf{Component} & \textbf{Description} & \textbf{Score} \\
\midrule
1 & Warm-Up        & Spiral review: distance from zero, distance between two numbers, and
                     shading a number line                                                    & \blank{1.2cm} \\
2 & Guided Notes   & Absolute value as distance; isolate then split; ``less th\emph{AND}'' vs.
                     ``great\emph{OR}''; the three notations --- ending in \emph{Guided
                     Practice} and \emph{Individual Practice}                                 & \blank{1.2cm} \\
3 & Homework       & Mixed practice: equations, a contrast pair, reading a number line, the
                     special cases, a model --- \emph{due next class}                         & \blank{1.2cm} \\
\midrule
  & \multicolumn{2}{r}{\textbf{Total}} & \blank{1.2cm} \\
\end{tabularx}
\end{tocbox}

\vspace{0.10in}

\begin{remindbox}
\small $|x-h|$ is the \textbf{distance} between $x$ and $h$, so it is never negative and two
different numbers sit the same distance out --- which is why $|X|=c$ (with $c>0$) has \emph{two}
answers: $X=c$ or $X=-c$. \textbf{Isolate the bars before you split.} For inequalities, \emph{less
th\textbf{AND}} ($<,\le$) means \emph{closer than} $c$, giving one interval $-c<X<c$; \emph{great\textbf{OR}}
($>,\ge$) means \emph{farther than} $c$, giving two rays $X<-c$ or $X>c$. Write the answer as a set,
as an interval (brackets include an endpoint, parentheses exclude it --- and $\infty$ always takes a
parenthesis), or as a shaded number line.
\end{remindbox}

\end{document}
